\documentclass[12pt]{article}
\usepackage{amsmath}
\usepackage[margin=1in]{geometry}
\usepackage{xcolor}
\usepackage{listings}

\lstset{
  language=Python,
  basicstyle=\ttfamily\small,
  keywordstyle=\bfseries,
  commentstyle=\itshape,
  stringstyle=\ttfamily,
  numbers=left,
  numberstyle=\tiny,
  stepnumber=1,
  numbersep=8pt,
  frame=single,
  breaklines=true,
  showstringspaces=false,
  tabsize=2
}

\title{Labwork 4: Neural Network}
\author{Trung Tran Minh - 2540024}
\date{}

\begin{document}
\maketitle

\section{Introduction}

In this lab, I implement a simple neural network from scratch.
The network is initialized from a text file. The program also supports two ways to initialize weights and biases:
random initialization and loading weights from another text file.

The network supports feedforward. I also test the network with the XOR problem.

\section{XOR Network Define}

The input file \texttt{xor\_network.txt} stores the network structure.
The file is:

\[
\begin{array}{c}
3 \\
2 \\
2 \\
1
\end{array}
\]

The first line means the number of layers is 3.
The next lines mean that the network has:

\[
2 \rightarrow 2 \rightarrow 1
\]

So the first layer has 2 input neurons, the hidden layer has 2 neurons, and the output layer has 1 neuron.

\newpage
\section{Implementation}

Below is the code that I implemented for the two classes \texttt{Layer} and \texttt{NeuralNetwork}.
The \texttt{Layer} class stores weights and biases of one layer.
The \texttt{NeuralNetwork} class creates all layers, runs feedforward, and can load weights from a text file.

\begin{lstlisting}[caption={Implementation of Layer and NeuralNetwork classes}]
class Layer:
  def __init__(self, input_size, output_size):
    self.input_size = input_size
    self.output_size = output_size
    self.weights = []
    self.biases = []

    for i in range(output_size):
      row = []

      for j in range(input_size):
        row.append(random.random())

      self.weights.append(row)
      self.biases.append(random.random())

class NeuralNetwork:
  def __init__(self, sizes):
    self.sizes = sizes
    self.layers = []

    for i in range(len(sizes) - 1):
      layer = Layer(sizes[i], sizes[i + 1])
      self.layers.append(layer)

  def sigmoid(self, x):
    if x < -100:
      return 0.0
    if x > 100:
      return 1.0

    return 1 / (1 + math.exp(-x))

  def feedforward(self, inputs):
    a = inputs

    for layer in self.layers:
      next_a = []

      for i in range(layer.output_size):
        z = layer.biases[i]

        for j in range(layer.input_size):
          z = z + layer.weights[i][j] * a[j]

        next_a.append(self.sigmoid(z))

      a = next_a

    return a

  def load_weights(self, file_path):
    with open(file_path, 'r') as file:
      lines = file.readlines()

    clean_lines = []

    for line in lines:
      line = line.strip()

      if line != '':
        clean_lines.append(line)

    k = 0

    for layer in self.layers:
      for i in range(layer.output_size):
        numbers = clean_lines[k].split(',')
        numbers = [float(x) for x in numbers]

        for j in range(layer.input_size):
          layer.weights[i][j] = numbers[j]

        layer.biases[i] = numbers[-1]
        k = k + 1

  def print_network(self):
    for layer_index in range(len(self.layers)):
      layer = self.layers[layer_index]
      print('Layer', layer_index + 1)

      for i in range(layer.output_size):
        print(' neuron', i + 1, 'weights =', layer.weights[i], 'bias =', layer.biases[i])
\end{lstlisting}

\section{mplement Feedforward}

For each neuron, I calculate:

\[
z = b + \sum_i w_i x_i
\]

Then I use sigmoid activation:

\[
a = \frac{1}{1 + e^{-z}}
\]

The output of one layer becomes the input of the next layer.
This continues until the output layer.

\section{XOR Network Experiment}

For the XOR experiment, I use this network:

\[
2 \rightarrow 2 \rightarrow 1
\]

The weight file \texttt{xor\_weights.txt} is:

\[
\begin{array}{c}
20, 20, -10 \\
-20, -20, 30 \\
20, 20, -30
\end{array}
\]

The first two rows are for the hidden layer.
The last row is for the output layer.
The last value in each row is the bias.

\section{Result}

I run the network with four XOR inputs.
The result is:

\[
\begin{array}{c|c|c|c}
x_1 & x_2 & output & predicted \\
\hline
0 & 0 & 0.000045 & 0 \\
0 & 1 & 0.999955 & 1 \\
1 & 0 & 0.999955 & 1 \\
1 & 1 & 0.000045 & 0
\end{array}
\]

The predicted XOR result is:

\[
0, 1, 1, 0
\]

This is the correct output for XOR.
\end{document}
